
%% ============================================================
%%  CHƯƠNG 1: TỔNG QUAN
%% ============================================================
\chapter{Tổng quan}

\section{Giới thiệu về Trung tâm An ninh mạng}
\label{sec:company}

Trung tâm An ninh mạng (Cyber Security Center) là đơn vị hoạt động trong lĩnh vực bảo mật thông tin và an toàn mạng, cung cấp các giải pháp phát hiện, ngăn chặn và ứng phó với các mối đe dọa mạng. Với đội ngũ chuyên gia giàu kinh nghiệm và hệ thống cơ sở hạ tầng kỹ thuật hiện đại, Trung tâm hướng đến việc nghiên cứu và phát triển các công nghệ tiên tiến trong lĩnh vực an ninh mạng, bao gồm hệ thống phát hiện xâm nhập (IDS), phân tích mã độc, săn tìm mối đe dọa (threat hunting) và ứng dụng trí tuệ nhân tạo trong bảo mật.

Các lĩnh vực hoạt động chính của Trung tâm bao gồm:
\begin{itemize}
    \item Nghiên cứu và phát triển các giải pháp an ninh mạng tiên tiến dựa trên trí tuệ nhân tạo và học máy.
    \item Xây dựng hệ thống giám sát an ninh mạng (SOC) và quy trình ứng phó sự cố.
    \item Phân tích mã độc và điều tra số (digital forensics).
    \item Đánh giá bảo mật hệ thống (penetration testing, vulnerability assessment).
    \item Đào tạo và nâng cao nhận thức về an ninh mạng cho tổ chức và cá nhân.
    \item Hợp tác nghiên cứu với các trường đại học, viện nghiên cứu trong và ngoài nước.
\end{itemize}

Môi trường thực tập tại Trung tâm tạo điều kiện cho sinh viên tiếp cận với các bài toán thực tế, hệ thống hạ tầng mạng quy mô lớn và những thách thức bảo mật đang nổi lên trong bối cảnh chuyển đổi số. Trong thời gian thực tập, nhóm tác giả đã được tham gia vào quá trình nghiên cứu và phát triển hệ thống phát hiện xâm nhập dựa trên học liên kết, một hướng tiếp cận mới nhằm giải quyết các vấn đề về quyền riêng tư dữ liệu, tính không đồng nhất của môi trường mạng và khả năng chống chịu trước các cuộc tấn công độc hại.

Quy trình thực tập tại Trung tâm được tổ chức theo các giai đoạn: (i) định hướng và làm quen với môi trường làm việc, (ii) nghiên cứu lý thuyết và tìm hiểu công nghệ nền tảng, (iii) thiết kế và phát triển giải pháp, (iv) thực nghiệm và đánh giá kết quả, (v) hoàn thiện báo cáo và trình bày kết quả. Trong suốt quá trình thực tập, nhóm tác giả được sự hướng dẫn trực tiếp từ các chuyên gia và cán bộ kỹ thuật của Trung tâm, đảm bảo định hướng nghiên cứu phù hợp với nhu cầu thực tế và xu hướng công nghệ mới nhất.

\section{Lý do chọn đề tài}
\label{sec:reason}

Hệ thống phát hiện xâm nhập (Intrusion Detection Systems -- IDS) là một thành phần quan trọng trong hạ tầng an ninh mạng hiện đại, cho phép xác định các hoạt động độc hại trong các mạng doanh nghiệp, đám mây, công nghiệp và Internet vạn vật (IoT). Các báo cáo gần đây cho thấy các cuộc tấn công nhắm vào IoT tiếp tục gia tăng \cite{sonicwall2024report}, trong khi các hệ thống bảo vệ nhà thông minh chặn hàng triệu sự kiện độc hại mỗi ngày \cite{netgear2024report}. Các mối đe dọa này bao gồm từ lây nhiễm mã độc, tuyển dụng botnet đến các chiến dịch tấn công từ chối dịch vụ phân tán (DDoS) quy mô lớn \cite{hackernews2026report}, làm nổi bật nhu cầu về các hệ thống phát hiện xâm nhập hiệu quả và có khả năng mở rộng.

Do đó, những tiến bộ gần đây trong Học máy (Machine Learning -- ML) và Học sâu (Deep Learning -- DL) đã cải thiện đáng kể khả năng phát hiện của IDS thông qua việc tự động học các mẫu lưu lượng phức tạp từ dữ liệu mạng quy mô lớn. Tuy nhiên, việc huấn luyện hiệu quả các IDS dựa trên ML thường yêu cầu thu thập dấu vết lưu lượng từ nhiều tổ chức, gây ra các vấn đề về quyền riêng tư, bảo mật và tuân thủ quy định. Học liên kết (Federated Learning -- FL) \cite{fedavg} đã nổi lên như một mô hình đầy hứa hẹn cho phép huấn luyện IDS cộng tác mà không chia sẻ trực tiếp dữ liệu lưu lượng riêng tư.

Mặc dù thiết kế của FL cải thiện rõ rệt tính bảo mật dữ liệu của các bên tham gia, việc triển khai thực tế FL-based IDS vẫn còn nhiều thách thức. Các mạng nơ-ron sâu hiện đại thường chứa hàng triệu tham số, dẫn đến chi phí truyền thông đáng kể trong quá trình trao đổi tham số lặp đi lặp lại. Hơn nữa, các nghiên cứu gần đây đã chỉ ra rằng các bản cập nhật mô hình có thể làm rò rỉ thông tin huấn luyện nhạy cảm thông qua các cuộc tấn công gradient inversion và reconstruction. Mặc dù tồn tại nhiều giải pháp cho vấn đề này như Secure Aggregation \cite{bonawitz2016secagg} và Homomorphic Encryption \cite{li2020deepfed}, những kỹ thuật này lại tạo ra chi phí tính toán và truyền thông đáng kể do ứng dụng các bước mã hóa lên trọng số mô hình.

Để giảm việc truyền tham số và hạn chế rủi ro lộ thông tin từ các bản cập nhật mô hình, các nghiên cứu gần đây đã ứng dụng kỹ thuật chắt lọc kiến thức (Knowledge Distillation) \cite{hinton2015kd} vào học máy liên kết, mở ra một lĩnh vực nghiên cứu học chắt lọc liên kết (Federated Distillation -- FD) \cite{li2024fdsurvey}. Thay vì trao đổi tham số mô hình, các phần tử client trong hệ thống FD trao đổi các biểu diễn tri thức trung gian như logits hoặc prototypes. Điều này cho phép các hệ thống FD điều chỉnh chi phí truyền thông đồng thời hỗ trợ liên kết giữa các kiến trúc mô hình không đồng nhất (model heterogeneity). Nhờ đó, FD đã trở thành một giải pháp thay thế hấp dẫn cho các môi trường hạn chế tài nguyên và không đồng nhất, đặc biệt là các thiết bị biên (edge devices), thiết bị nhúng (IoT) và nút mạng.

Tuy nhiên, hầu hết các giải pháp FD hiện tại vẫn phụ thuộc vào một máy chủ tổng hợp tập trung. Trong các triển khai quy mô lớn, sự ổn định của hệ thống học liên kết phải phụ thuộc vào đúng một phần tử duy nhất là máy chủ. Nếu máy chủ này gặp sự cố do tấn công hoặc lỗi hạ tầng, toàn bộ quá trình học liên kết bị buộc phải ngừng trệ. Do đó, học liên kết phân tán (Decentralized Federated Learning -- DFL) giải quyết hạn chế này bằng cách phân phối trách nhiệm điều phối giữa các clients tham gia và loại bỏ sự phụ thuộc vào máy chủ vĩnh viễn \cite{braintorrent}.

Ngoài các vấn đề về băng thông truyền tải và điều phối quy trình học liên kết, các hệ thống liên kết vẫn dễ bị tổn thương trước các tác nhân Byzantine cố tình đầu độc dữ liệu cục bộ (data poisoning), đầu độc cập nhật mô hình (model poisoning) hoặc đầu độc biểu diễn trung gian (logit poisoning) \cite{tolpegin2020targetedflip, fang2020attacks, roux2025byzantine}. Các cuộc tấn công gần đây đã tiến hóa từ việc tạo ra các giá trị cực đại dễ phát hiện đến việc tạo ra các bản cập nhật độc hại gần giống với hành vi bất thường trong ngữ cảnh bất đồng dữ liệu Non-IID, khiến việc phát hiện ngày càng khó khăn \cite{roux2025byzantine}. Các phương pháp phòng thủ tấn công Byzantine hiện có hầu hết tập trung vào FL trong không gian tham số, trong khi các phương pháp phòng thủ có thể ứng dụng cho không gian biểu diễn (như logit) cho các hệ thống FD vẫn còn tương đối mới và ít được nghiên cứu hơn.

Trong giai đoạn thực tập, em và thành viên trong nhóm nghiên cứu đã xây dựng, thực nghiệm và đề xuất DFD-IDS, một cơ chế chắt lọc kiến thức liên kết phân tán có khả năng chống chịu đầu độc Byzantine, với cơ chế tổng hợp logits dựa trên energy. DFD-IDS kết hợp Chắt lọc kiến thức dựa trên bằng chứng (Evidential Knowledge Distillation -- EKD) để bảo toàn độ bất định trong phân phối dự đoán của mô hình, kỹ thuật tiết chế đóng góp dựa trên energy (Energy-based Vote Abstention -- EVA) để lọc các dự đoán mẫu công khai thiếu tự tin, cơ chế lọc kết hợp RobustFilter \cite{diakonikolas2017robustfilter} và ước lượng xác suất phân phối chuẩn của Blom \cite{blom} mới có khả năng chống chịu Byzantine hiệu quả cao. Thêm vào đó, thiết kế lựa chọn phần tử điều phối cho môi trường phân tán được thực hiện dựa trên điểm số lọc của từng client. Bằng cách kết hợp các thành phần này, giải pháp đề xuất cho phép chia sẻ tri thức hiệu quả giữa các clients không đồng nhất về cả kiến trúc mô hình và phân bố dữ liệu, đồng thời duy trì khả năng chống chịu trước các cuộc tấn công đầu độc Byzantine trong môi trường phân tán.

Các đóng góp chính của đề tài bao gồm:
\begin{itemize}
    \item Đề xuất DFD-IDS, một hệ thống học chắt lọc kiến thức liên kết phi tập trung cho phát hiện xâm nhập phân tán có khả năng phòng chống tấn công Byzantine.
    \item Giới thiệu cơ chế Energy-based Vote Abstention (EVA) giúp thích ứng phát hiện ngoài phân phối dựa trên năng lượng vào việc chọn lọc dự đoán mẫu công khai trong môi trường liên kết, mà không cần lệ thuộc trên một mô hình phụ để phân biệt.
    \item Phát triển Bộ lọc Blom Robust Filter giúp quá trình tổng hợp logits chống được tấn công Byzantine trong môi trường clients không đồng nhất.
    \item Đánh giá DFD-IDS trước các kỹ thuật tấn công Byzantine hiện đại và chứng minh khả năng chống chịu vượt trội so với các phương pháp cơ sở hiện có.
\end{itemize}

\section{Mục tiêu, đối tượng và phạm vi}
\label{sec:objectives}

\subsection{Mục tiêu nghiên cứu}
Mục tiêu chính của đề tài là nghiên cứu, thiết kế và hiện thực một hệ thống phát hiện xâm nhập mạng dựa trên học liên kết chắt lọc có khả năng phân tán, chống chịu trước các tấn công đầu độc (poisoning attacks) và hỗ trợ kiến trúc mô hình không đồng nhất (heterogeneous architectures). Cụ thể:

\begin{itemize}
    \item Nghiên cứu tổng quan về các hệ thống IDS dựa trên học máy và học liên kết.
    \item Tìm hiểu các phương pháp Federated Distillation và khả năng ứng dụng trong bảo mật mạng.
    \item Phân tích các dạng tấn công độc hại trong học liên kết và cơ chế phòng thủ tương ứng.
    \item Đề xuất và hiện thực khung giải pháp DFD-IDS với các thành phần chống chịu Byzantine.
    \item Đánh giá hiệu quả của giải pháp thông qua các kịch bản thực nghiệm.
\end{itemize}

\subsection{Đối tượng nghiên cứu}
Đối tượng nghiên cứu của đề tài bao gồm:
\begin{itemize}
    \item Hệ thống phát hiện xâm nhập mạng (Network-based IDS).
    \item Học liên kết tập trung và phân tán (Centralized \& Decentralized Federated Learning).
    \item Các phương pháp chắt lọc kiến thức liên kết (Federated Distillation).
    \item Các kỹ thuật phát hiện và chống chịu tấn công độc hại trong môi trường liên kết.
    \item Mô hình học sâu (Deep Neural Networks, GRU, CNN) cho dữ liệu mạng.
\end{itemize}

\subsection{Phạm vi nghiên cứu}
Đề tài tập trung vào:
\begin{itemize}
    \item Dữ liệu mạng từ các bộ dữ liệu chuẩn như CSE-CIC-IDS2018 và các tập dữ liệu liên quan.
    \item Môi trường thực nghiệm trong phòng lab với các máy chủ được cấu hình phù hợp.
    \item Đánh giá trên các chỉ số: Accuracy, Macro precision, Macro recall, Macro F1-Score và khả năng chống chịu tấn công.
\end{itemize}

\section{Những điểm mới của đề tài}
\label{sec:novelty}

Đề tài nghiên cứu có những điểm mới sau đây so với các công trình hiện có:

\begin{enumerate}
    \item \textbf{Kiến trúc phân tán hoàn toàn}: Không giống như các phương pháp FL-IDS hiện tại chủ yếu dựa trên server trung tâm, DFD-IDS loại bỏ hoàn toàn sự phụ thuộc vào một máy chủ trung tâm thông qua cơ chế chọn phần tử điều phối từ kết quả tổng hợp logit. Điều này giúp loại bỏ lệ thuộc vào một phần tử duy nhất, qua đó cải thiện tính ổn định hệ thống trong vận hành thực tế.
    \item \textbf{Energy-based Vote Abstention (EVA)}: Đề xuất áp dụng cơ chế phát hiện out-of-distribution dựa trên năng lượng (Energy-based OOD) vào bài toán chọn lọc dự đoán nhãn công khai trong Federated Distillation. EVA tiết kiệm chi phí hơn nhiều so với phương pháp dựa vào mô hình phụ của SSFL-IDS và không yêu cầu tạo sinh một bộ dữ liệu riêng cho mô hình này.
    \item \textbf{Blom Robust Filter}: Đề xuất một phương pháp tổng hợp mới dựa trên xác suất phân phối Blom, khắc phục hạn chế từ phương pháp RobustFilter phải lệ thuộc vào một ngưỡng cố định như đã ứng dụng trong Cronus \cite{chang2019cronus}. Phương pháp này có khả năng thích ứng với số lượng clients và số chiều logits khác nhau.
    \item \textbf{Kết hợp EKD và FL}: Ứng dụng kỹ thuật chắt lọc hiện đại là Evidential Knowledge Distillation để cho phép các kiến trúc mô hình khác nhau tham gia học liên kết.
\end{enumerate}

\section{Phương pháp nghiên cứu}
\label{sec:methodology}

Phương pháp nghiên cứu được áp dụng trong đề tài bao gồm:
\begin{itemize}
    \item \textbf{Nghiên cứu lý thuyết}: Thu thập, phân tích và tổng hợp các tài liệu khoa học liên quan đến IDS, Federated Learning, Federated Distillation, Byzantine-robust aggregation và các lĩnh vực liên quan.
    \item \text{Nghiên cứu thực nghiệm}: Thiết kế và hiện thực các mô hình, tiến hành thực nghiệm trên các bộ dữ liệu chuẩn, so sánh kết quả với các phương pháp nền tảng.
    \item \textbf{Đánh giá định lượng}: Sử dụng các chỉ số đánh giá khách quan để đo lường hiệu suất của hệ thống.
    \item \textbf{Phân tích so sánh}: So sánh phương pháp đề xuất với các phương pháp hiện có (FedAvg, FedProto, FLAME, SSFL-IDS và FedKD-IDS) trên cùng điều kiện thực nghiệm.
\end{itemize}

\section{Kế hoạch thực tập chi tiết}
\label{sec:plan}

Quá trình thực tập tại Trung tâm An ninh mạng được chia thành 8 tuần với các mục tiêu cụ thể cho từng giai đoạn. Bảng~\ref{tab:plan_summary} tóm tắt kế hoạch thực tập tổng thể.

\begin{table}[H]
\centering
\caption{Kế hoạch thực tập tổng thể}
\label{tab:plan_summary}
\begin{tabular}{|c|c|c|}
\hline
\textbf{Giai đoạn} & \textbf{Tuần} & \textbf{Mục tiêu} \\
\hline
Định hướng & 1-2 & Làm quen môi trường, nghiên cứu tổng quan \\
\hline
Nghiên cứu lý thuyết & 3-4 & Tìm hiểu FL, FD, Byzantine attacks, defenses \\
\hline
Phát triển giải pháp & 5-6 & Thiết kế và hiện thực DFD-IDS \\
\hline
Thực nghiệm \& Đánh giá & 7 & Tiến hành thực nghiệm, thu thập kết quả \\
\hline
Hoàn thiện & 8 & Phân tích kết quả, viết báo cáo \\
\hline
\end{tabular}
\end{table}

\section{Cấu trúc báo cáo}
\label{sec:structure}

Báo cáo thực tập được tổ chức thành 5 chương:

\begin{itemize}
    \item \textbf{Chương 1 -- Tổng quan}: Giới thiệu về Trung tâm An ninh mạng, bối cảnh nghiên cứu, lý do chọn đề tài, mục tiêu, đối tượng, phạm vi và phương pháp nghiên cứu.
    \item \textbf{Chương 2 -- Cơ sở lý thuyết}: Trình bày các kiến thức nền tảng hệ thống IDS, học máy liên kết, chắt lọc kiến thức, học chắt lọc liên kết, tấn công Byzantine và các kỹ thuật phòng thủ tương ứng.
    \item \textbf{Chương 3 -- Phương pháp thực hiện}: Mô tả chi tiết khung giải pháp DFD-IDS được đề xuất, bao gồm kiến trúc tổng thể, cơ chế Energy-based Vote Abstention, Blom Robust Filter, cơ chế chọn server phân tán và Evidential Knowledge Distillation.
    \item \textbf{Chương 4 -- Thực nghiệm và đánh giá}: Trình bày môi trường thực nghiệm, tập dữ liệu, các kịch bản thử nghiệm, kết quả đạt được và bảng so sánh với các phương pháp khác.
    \item \textbf{Chương 5 -- Kết luận và hướng phát triển}: Tổng kết các kết quả đạt được, đánh giá ưu nhược điểm của giải pháp và đề xuất các hướng phát triển trong tương lai.
\end{itemize}

\pagebreak
